\chapter{引言}\label{chap:introduction}

\section{研究背景}

考虑如下场景：分拣传送带上，立方体零件持续平移和旋转，机械臂需逐一抓取并将其搬运至指定目标区域。与”传送带停止—抓取—启动”的间歇模式不同，动态抓取要求机械臂在目标运动过程中依次完成两项任务：其一，接近目标并实现相对位姿对齐；其二，在恰当时机闭合夹爪，保持抓取并搬运至指定区域\citep{marturi2019dynamic,zhang2022robotic,newbury2023deep}。机械臂需要跟踪当前目标位置、预测目标短期运动，同时可靠地判断当前任务阶段——夹爪闭合过早会导致对齐不足，闭合过晚则可能错过可达窗口。

现有的动态抓取方法大致分为两类。基于模型的方法（如FSM-MPC）能够显式处理位置、姿态、平滑性约束和执行约束，在状态观测准确时结构清晰、可解释性强\citep{rawlings2017model,nubert2020safe}。然而，目标位姿观测存在噪声、离群点或时延时，手工阶段规则和参考轨迹构造会将感知误差直接传递至控制输出，导致阶段判定和跟踪参考系统性地偏离真实状态。另一方面，模仿学习从专家示教中学习动态追踪、相对位姿对齐和夹爪闭合时机\citep{argall2009survey}，但离线策略在闭环执行时会受到分布偏移和误差累积的影响\citep{ross2010efficient,ross2011reduction}。

上述两类方法各有优劣，但更根本的问题在于：策略架构、观测噪声、状态估计策略和执行器选择等因素相互纠缠，闭环成功率的差异究竟源自何处？已有工作中，这些因素通常随输入输出表示和执行协议的不同而变化，彼此难以解耦，因此各因素的独立贡献无法被单独衡量。

本文构建了统一的实验框架，系统分析了影响动态抓取闭环成功率的关键因素。该框架的核心是统一的目标中心动作块模仿学习接口，使所有策略共享相同的观测--动作表示、运行时阶段约束和MPC跟踪执行器。基于该框架，本文以大规模专家示教为基础，比较了BC、BC-RNN、ACT和DP四类策略。

\section{相关工作}
\label{sec:related-work}

\subsection{动态目标抓取}

动态目标抓取要求机器人在目标持续运动过程中完成相对位姿对齐和接触决策，难度远高于静态抓取。现有抓取综述通常涵盖解析抓取、采样规划、学习抓取和闭环视觉伺服等路线，探讨实现鲁棒操作的不同技术路径\citep{zhang2022robotic,newbury2023deep}。目标运动时，机器人还需估计其运动趋势，在有限时间内确定接近路径、抓取面、夹爪朝向和闭合时机\citep{marturi2019dynamic,alshanoon2022robotic}。这意味着动态抓取需同时完成感知、预测、阶段判断与执行跟踪，而非仅依赖单帧抓取姿态估计。

传统动态抓取方法通常依赖几何预测、视觉跟踪、轨迹优化或显式状态机。状态观测准确且运动模型可用时，这类方法解释性强、稳定性好；但目标姿态、速度或感知噪声变化较大时，手工阶段规则和参考轨迹构造可能变得脆弱。

本文以翻滚场景（立方体初始姿态均匀随机、绕空间轴旋转）中的立方体抓取为核心任务，系统分析目标运动、姿态对齐和抓取阶段时机对闭环成功率的影响。专家FSM-MPC控制器同时作为示教生成器和对照基线：其一，检验基于模型的控制器在噪声观测条件下是否必须依赖显式状态估计；其二，检验学习策略能否通过隐式利用时间冗余形成另一种噪声鲁棒机制。

现有动态抓取研究往往使用不同的评估标准、执行器和观测表示，因此很难判断不同工作之间的性能差异，究竟是策略设计所致还是实验设置使然。本文以固定评估标准、共享观测--动作接口和受控消融设计，搭建了从感知到执行的完整比较框架。

完整的动态目标抓取包含两个相互关联但可分离的子问题：抓取点选择（在物体表面确定抓取位置和夹爪朝向）和轨迹跟踪与时机控制（在目标运动过程中完成接近、对准和闭合）。抓取合成研究从解析方法演进至基于学习的方法，已成为独立的技术方向\citep{bohg2014grasp,newbury2023deep,morrison2018ggcnn,mahler2017dexnet}；拦截抓取问题的早期工作则从模型预测和控制角度处理运动目标的跟踪与抓取\citep{allen1993automated}。本文聚焦于后者——在给定抓取点（立方体中心，从上方抓取）的条件下，研究目标中心表示、时间动作块结构和噪声鲁棒性策略对闭环成功率的影响。本文有意将研究范围限定于该设定：固定抓取点后，实验比较可以更清晰地分离表示、时序结构和执行器因素的效应，而不受抓取点选择变化的干扰。这一设计选择的局限性与泛化范围将在第五章进一步讨论。

\subsection{模仿学习与闭环分布偏移}

模仿学习通过专家示教学习从观测到动作的映射，是机器人操作中常用的数据驱动策略学习方法\citep{argall2009survey}。行为克隆将专家轨迹分解为监督学习样本，这一思想可追溯到端到端驾驶策略ALVINN\citep{pomerleau1989alvinn}。但离线训练策略的动作会在闭环执行中改变后续观测分布；一旦进入训练数据未能覆盖的分布外状态，误差可能随时间累积并导致任务失败\citep{ross2010efficient,ross2011reduction}。动态抓取尤其容易出现这类问题，因为轻微的相对位姿误差会改变后续可抓取窗口和夹爪闭合时机。

交互式数据聚合和更大规模示教可以缓解分布偏移\citep{ross2011reduction,mandlekar2021matters}，但本文使用算法生成的离线专家示教，因此更聚焦于结构设计对闭环执行的影响。本文将单步BC作为基础参照。若单步回归在翻滚场景中即可成功，则说明目标中心表示和监督信号已足够应对该任务；此时闭环分布偏移的缓解应归因于历史条件窗口、动作块重规划和阶段约束的组合，而非数据聚合。

现有离线模仿学习研究在策略比较中，通常由于观测表示、动作空间和执行协议不一致，难以分离架构效应与实验设置效应。举例来说，观测表示的变化给同一策略带来的性能波动，可能比同一表示下不同策略之间的差异更大。本文要求所有策略共享相同的观测--动作接口和运行时结构，将表示效应、架构效应和执行器效应分离为独立消融维度，进而量化各因素的相对贡献。

\subsection{时间建模与动作分块}

时间建模是动态操作任务中的关键设计。循环网络通过隐藏状态整合历史观测，在部分可观测、多阶段操作或示教质量变化的任务中构成了强基线\citep{mandlekar2021matters}。BC-RNN通常使用GRU或LSTM编码历史信息，让策略能够利用目标速度、相对相位和近期执行误差等时序线索。在动态抓取任务中，这些线索直接关系到抓取面选择、对齐阶段持续时间和夹爪闭合时机。

动作分块进一步将策略输出从单步动作扩展为短时动作序列。ACT将机器人操作策略建模为条件动作序列生成问题，利用Transformer编码和条件序列生成提高闭环动作的时序一致性\citep{zhao2023learning,vaswani2017attention}。在本文框架内，BC、BC-RNN、ACT和DP都输出32步任务空间动作块，并在运行时每16步重规划一次，因此结果差异并非源自动作块接口的有无。本文关注的核心问题是：在统一动作块接口和翻滚场景条件下，循环结构、Transformer-CVAE结构和扩散去噪结构的相对性能差异。为此，本文在该场景中开展历史长度消融，检验更长历史窗口是否对循环和Transformer策略持续有益，抑或存在收益边界。

现有时间策略比较中有一个问题尚未解决：ACT的原始论文报告其性能优于BC-RNN，但该工作使用了不同的观测表示和执行设置；因此无法确定ACT的增益来自Transformer-CVAE架构、动作块结构还是表示选择。本文基于统一表示和接口，直接比较不同时间机制（GRU循环、Transformer注意力、扩散去噪）的闭环表现。这一设计使时间结构成为唯一变量，从而分离时间机制的独立效应。

\subsection{生成式动作策略与扩散策略}

机器人操作中的动作分布通常呈现多模态特性：在相同观测下，可能存在多种可行的接近方向、抓取面或闭合时机。混合密度模型、能量模型和生成式模型都曾用于刻画这类一对多映射\citep{bishop1994mixture,florence2022implicit}。扩散模型通过逐步加噪和反向去噪学习数据分布\citep{ho2020denoising,song2021denoising}，扩散策略（DP）则将这一思想应用于机器人动作序列生成，在视觉操作任务中展现了较强的轨迹生成能力\citep{chi2023diffusion}。

生成式动作建模本身并不确保更高的闭环成功率。DP的闭环表现受多种设计因素影响，包括状态表示、动作参数化、模型容量、采样步数和跟踪执行器等。本文在低维目标中心任务空间接口中比较DP与确定性动作块策略，同时检验DP在不同容量和世界坐标表示下的表现。此外，本文还检验了验证噪声预测损失是否可作为闭环成功率的替代指标。通过比较验证损失与闭环成功率的变化趋势，本文将生成式建模能力与其在闭环控制中的实际收益加以区分——更低的验证损失并不保证更高的闭环成功率。

现有扩散策略研究通常同时报告验证损失和闭环成功率，但尚未系统检验二者之间的相关性——更低的验证损失是否对应更高的闭环成功率。原始Diffusion Policy研究虽报告了不同容量和表示配置下的结果，但并未围绕验证损失与闭环成功率的关系开展受控分析。此外，扩散模型在机器人操作中的应用并不限于动作分块去噪——Diffuser等方法在轨迹级别实施扩散规划\citep{janner2022diffuser}，其设计和接口选择不同于DP。本文的DP实验基于固定推理步数的条件扩散去噪实现，因此相关结论仅限于该特定配置下的表现，不推广至其他扩散策略变体。本文在同一框架内系统改变DP容量（1.5M至18M）和表示（目标中心与世界坐标），直接检验验证损失与闭环成功率之间是否存在单调关系。

\subsection{目标中心表示与相对动作}

机器人操作中的坐标表示显著影响学习难度。若直接在世界坐标系中学习末端位姿参考，同一种抓取意图会随目标位置和姿态变化而产生数值差异很大的动作目标；而在目标或物体坐标系中表示动作，策略可以直接学习”相对目标如何接近、对齐和闭合”。类似的相对轨迹或物体中心表示已被用来提高机器人操作策略对空间变化的泛化能力\citep{zeng2021transporter,chi2024universal}。

本文将历史观测、未来目标条件和动作标签均表示于目标中心坐标系。通过系统比较目标中心表示与世界坐标表示（各含完整和紧凑两种变体），直接量化表示选择对闭环成功率的影响。关键问题是：表示选择与架构选择，究竟哪一个对闭环成功率的影响更大？如果目标中心与世界坐标之间的成功率差距超过BC与ACT之间的差距，则说明在动态抓取任务中，表示设计比模型选择更为根本。

已有研究分别证实了目标中心表示的有效性和特定架构的优越性，但尚未在同一实验中同时量化两者的影响程度并直接比较。本文通过交叉表示类型与策略类型的全因子实验设计补充了这一缺失：在统一评估协议下同时报告表示效应与架构效应及其95\% Wilson置信区间，读者可以直接比较两类设计选择的相对重要性，无需依赖不同实验设置的结果来推断。

\subsection{混合控制与跟踪执行器}

学习策略直接输出关节命令可以减少对手工结构的依赖，但接触约束、相位逻辑和低层跟踪也因此完全交由数据驱动模型来习得。分层模仿学习和学习规划器方法通常将任务意图、局部轨迹和低层执行分开建模\citep{lynch2020learning,mandlekar2020iris,wang2023mimicplay}。本文借鉴类似思想，但低层执行仍由结构化跟踪执行器负责：学习策略预测任务空间动作块，FSM约束阶段推进，MPC跟踪执行器负责跟踪任务空间参考。这一设计确保策略比较主要反映动作块生成质量，而非执行器的差异。

现有混合控制框架通常将学习策略与固定执行器绑定评估，很少在同一实验中控制执行器差异。因此，不同策略采用不同执行器时，性能差异同时反映策略设计和执行器能力，因而难以归因。本文主实验统一使用MPC跟踪执行器，与专家结构保持一致，从而将执行器差异从策略架构比较中剔除。综上，本文后续章节基于统一的观测--动作接口和运行时结构，将表示效应、架构效应和执行器效应分离为独立消融维度，以系统回答前述文献中尚未解决的问题。

\section{研究问题}

本文以翻滚环境为主要实验平台：立方体初始姿态均匀随机，绕空间轴以中等角速度旋转并同时平移。本文中”翻滚”指三维姿态绕非固定竖直轴变化，而非立方体与地面接触滚动。同时，本文在较简单的定轴自旋场景（立方体初始姿态固定、绕$z$轴自旋）中进行对照评估，以确认翻滚场景中的性能下降源于任务难度，而非评估协议或场景设定差异（详见第四章\S\ref{sec:main-clean}）。

基于上述设置，本文聚焦三个核心问题：

\begin{enumerate}
    \item 表示选择：目标中心表示是否优于世界坐标表示——更重要的是，这一影响是否超过架构选择？已有工作通常预设某一坐标表示，但很少将其作为独立变量进行消融研究。本文在目标中心表示（完整和紧凑）与世界坐标表示（完整和紧凑）之间直接比较，检验表示选择对BC和ACT的影响程度。
    \item 时间结构：翻滚场景是否需要超越单步回归的时间结构，且历史窗口长度如何影响循环和注意力策略？该场景的难点在于随机初始姿态和空间轴旋转使抓取面选择和姿态对齐变得复杂。本文以单步BC检验时序信息的必要性，以BC-RNN和ACT检验循环编码和Transformer查询能否弥补这一缺陷，并通过历史长度消融分析时间上下文的收益边界。
    \item 噪声鲁棒性：应通过显式状态滤波实现，还是依赖策略自身的隐式时间冗余利用？本文分别检验专家FSM-MPC控制器和学习策略在引入卡尔曼滤波后的表现，以区分显式状态估计和隐式时间冗余利用在动态抓取中的互补角色。
\end{enumerate}

\section{方法概述与主要贡献}

本文采用由学习规划器和结构化跟踪执行器组成的混合控制框架，如图\ref{fig:intro-overview}所示。由有限状态机和MPC组成的专家控制器首先生成成功示教，学习策略随后基于目标中心历史观测和短时未来目标条件预测32步任务空间相对动作块。在线执行时，系统每16步重新规划一次，并由单调相位状态机与MPC跟踪执行器共同完成阶段推进和轨迹跟踪。

本文的主要贡献有三：

\begin{enumerate}
    \item 构建动态立方体抓取实验框架。以翻滚场景为主要设置、定轴自旋场景为验证对照，以闭环成功率为核心评估指标（每条件750回合），覆盖无噪观测和强噪声观测两种条件。
    \item 提出统一的目标中心动作块模仿学习接口。该接口使BC、BC-RNN、ACT和DP共享输入输出表示、阶段约束和评估标准，从而将架构差异与表示差异和执行器差异解耦。
    \item 在以FSM-MPC为专家的实验框架内，系统实验明确了决定动态抓取闭环成功率的三个关键因素。其一，目标中心表示带来的提升最为稳定——在翻滚场景中，表示选择造成的成功率差距超过了架构选择的差距。其二，时间结构策略（BC-RNN、ACT）显著提高了翻滚场景中的闭环成功率。其三，噪声鲁棒性存在两条互补路径：显式状态滤波对模型控制器不可或缺，隐式时间冗余利用对学习策略已足够。
\end{enumerate}

\begin{figure}[!htbp]
    \centering
    \includegraphics[width=0.95\textwidth]{fig_intro_overview}
    \caption{目标中心动作块模仿学习框架。专家示教由FSM+MPC控制器生成，学习策略基于目标中心表示预测任务空间动作块，运行时通过单调相位状态机和MPC跟踪执行器实现闭环控制。}
    \label{fig:intro-overview}
\end{figure}

\section{章节安排}

本文其余章节安排如下。

第二章介绍动态抓取任务设置与专家控制器：涵盖仿真平台与评估协议，定轴自旋和翻滚两个实验场景的定义与设计意图，FSM-MPC专家架构与五阶段逻辑，以及专家示教数据集的采集与筛选。

第三章介绍统一模仿学习策略与实验设计：涵盖目标中心表示与动作块构造，BC、BC-RNN、ACT和DP四类策略的统一接口与训练协议，以及闭环执行与评估标准。

第四章报告实验结果：以定轴自旋和翻滚场景的双场景对比为主线，依次分析无噪观测下的跨场景策略比较、观测表示与历史结构消融、噪声鲁棒性与状态估计、时间集成消融，以及数据规模与模型容量边界。

第五章总结全文，讨论研究局限并展望未来工作。
